% GENERATED FILE. Edit canonical lessons, then run npm run generate:books.
% canonical-source-hash: fnv1a64:5f1d744051b95d58
% canonical-lessons: SA-C118-gulphah, SA-C118-angulih, SA-C118-kaksah, SA-C118-asandi, SA-C118-phalakam

\chapter{Ankle, Finger, Room, Chair, Board}
\label{ch:sa-ankle-finger-room-chair-board}
\hlchaptermodality{118}

\begin{chapteropening}
\textbf{By the end of this chapter:} \emph{I can say an ankle, a finger, a room, a chair and a board.}

Complete the last lesson of chapter 118: i can say an ankle, a finger, a room, a chair and a board.
\end{chapteropening}

\section[\texorpdfstring{\sk{गुल्फः} (gulphaḥ)}{गुल्फः (gulphaḥ)}]{\sk{गुल्फः} (gulphaḥ) — an ankle}

\label{lesson:SA-C118-gulphah}



\begin{quote}
Before the new one: say the Sanskrit for an eyebrow, then the Sanskrit for the skin.
\end{quote}

\subsection*{You'll want to know: \sk{गुल्फः}}
\textbf{\sk{गुल्फः}} — \emph{gulphaḥ} — \textquotedblleft{}an ankle\textquotedblright{}.

\begin{grammarlens}[title={What you've built}]
One more word for this chapter.
\end{grammarlens}

\subsection*{Your turn}
\begin{itemize}
\raggedright
  \item \emph{Say it:} \emph{gulphaḥ}
  \item \emph{Say it:} \emph{gulphaḥ}, once more
  \item \emph{Say it:} say \emph{gulphaḥ}
  \item \emph{Recall:} say the Sanskrit for an eyebrow, then the Sanskrit for the skin, then say \emph{gulphaḥ} again
\end{itemize}

\begin{tcolorbox}[breakable,colback=teal!4,colframe=teal!35!black,title={Before you move on}]
What does \sk{गुल्फः} mean? (\textquotedblleft{}An ankle\textquotedblright{}.) Say it once more.
\end{tcolorbox}
\section[\texorpdfstring{\sk{अंगुलिः} (aṅguliḥ)}{अंगुलिः (aṅguliḥ)}]{\sk{अंगुलिः} (aṅguliḥ) — a finger}

\label{lesson:SA-C118-angulih}



\begin{quote}
Before the new one: say the Sanskrit for the skin, then the Sanskrit for an ankle.
\end{quote}

\subsection*{You'll want to know: \sk{अंगुलिः}}
\textbf{\sk{अंगुलिः}} — \emph{aṅguliḥ} — \textquotedblleft{}a finger\textquotedblright{}.

\begin{grammarlens}[title={What you've built}]
One more word for this chapter.
\end{grammarlens}

\subsection*{Your turn}
\begin{itemize}
\raggedright
  \item \emph{Say it:} \emph{aṅguliḥ}
  \item \emph{Say it:} \emph{aṅguliḥ}, once more
  \item \emph{Say it:} say \emph{aṅguliḥ}
  \item \emph{Recall:} say the Sanskrit for the skin, then the Sanskrit for an ankle, then say \emph{aṅguliḥ} again
\end{itemize}

\begin{tcolorbox}[breakable,colback=teal!4,colframe=teal!35!black,title={Before you move on}]
What does \sk{अंगुलिः} mean? (\textquotedblleft{}A finger\textquotedblright{}.) Say it once more.
\end{tcolorbox}
\section[\texorpdfstring{\sk{कक्षः} (kakṣaḥ)}{कक्षः (kakṣaḥ)}]{\sk{कक्षः} (kakṣaḥ) — a room}

\label{lesson:SA-C118-kaksah}



\begin{quote}
Before the new one: say the Sanskrit for an ankle, then the Sanskrit for a finger.
\end{quote}

\subsection*{You'll want to know: \sk{कक्षः}}
\textbf{\sk{कक्षः}} — \emph{kakṣaḥ} — \textquotedblleft{}a room\textquotedblright{}.

\begin{grammarlens}[title={What you've built}]
One more word for this chapter.
\end{grammarlens}

\subsection*{Your turn}
\begin{itemize}
\raggedright
  \item \emph{Say it:} \emph{kakṣaḥ}
  \item \emph{Say it:} \emph{kakṣaḥ}, once more
  \item \emph{Say it:} say \emph{kakṣaḥ}
  \item \emph{Recall:} say the Sanskrit for an ankle, then the Sanskrit for a finger, then say \emph{kakṣaḥ} again
\end{itemize}

\begin{tcolorbox}[breakable,colback=teal!4,colframe=teal!35!black,title={Before you move on}]
What does \sk{कक्षः} mean? (\textquotedblleft{}A room\textquotedblright{}.) Say it once more.
\end{tcolorbox}
\section[\texorpdfstring{\sk{आसन्दी} (āsandī)}{आसन्दी (āsandī)}]{\sk{आसन्दी} (āsandī) — a chair}

\label{lesson:SA-C118-asandi}



\begin{quote}
Before the new one: say the Sanskrit for a finger, then the Sanskrit for a room.
\end{quote}

\subsection*{You'll want to know: \sk{आसन्दी}}
\textbf{\sk{आसन्दी}} — \emph{āsandī} — \textquotedblleft{}a chair\textquotedblright{}.

\begin{grammarlens}[title={What you've built}]
One more word for this chapter.
\end{grammarlens}

\subsection*{Your turn}
\begin{itemize}
\raggedright
  \item \emph{Say it:} \emph{āsandī}
  \item \emph{Say it:} \emph{āsandī}, once more
  \item \emph{Say it:} say \emph{āsandī}
  \item \emph{Recall:} say the Sanskrit for a finger, then the Sanskrit for a room, then say \emph{āsandī} again
\end{itemize}

\begin{tcolorbox}[breakable,colback=teal!4,colframe=teal!35!black,title={Before you move on}]
What does \sk{आसन्दी} mean? (\textquotedblleft{}A chair\textquotedblright{}.) Say it once more.
\end{tcolorbox}
\section[\texorpdfstring{\sk{फलकम्} (phalakam)}{फलकम् (phalakam)}]{\sk{फलकम्} (phalakam) — a board, a plank}

\label{lesson:SA-C118-phalakam}



\begin{quote}
Before the new one: say the Sanskrit for a room, then the Sanskrit for a chair.
\end{quote}

\subsection*{You'll want to know: \sk{फलकम्}}
\textbf{\sk{फलकम्}} — \emph{phalakam} — \textquotedblleft{}a board, a plank\textquotedblright{}.

\begin{grammarlens}[title={What you've built}]
One more word for this chapter.
\end{grammarlens}

\subsection*{Your turn}
\begin{itemize}
\raggedright
  \item \emph{Say it:} \emph{phalakam}
  \item \emph{Say it:} \emph{phalakam}, once more
  \item \emph{Say it:} say \emph{phalakam}
  \item \emph{Recall:} say the Sanskrit for a room, then the Sanskrit for a chair, then say \emph{phalakam} again
\end{itemize}

\begin{tcolorbox}[breakable,colback=teal!4,colframe=teal!35!black,title={Before you move on}]
What does \sk{फलकम्} mean? (\textquotedblleft{}A board, a plank\textquotedblright{}.) Say it once more.
\end{tcolorbox}
